\documentclass[11pt]{article}
\usepackage[a4paper,margin=28mm]{geometry}
\usepackage{amsmath,amssymb,amsthm,mathtools,hyperref}
\newcommand{\ddc}{dd^c}
\newtheorem{theorem}{Theorem}
\newtheorem{lemma}[theorem]{Lemma}
\theoremstyle{remark}\newtheorem{remark}[theorem]{Remark}
\title{A projective radial counterexample to\newline positive-volume bigness for nonclosed forms}
\author{}
\date{}
\begin{document}
\maketitle
\begin{abstract}
For every $n\geq3$ we construct a smooth semipositive nonclosed form on
$\mathbb P^{n-1}\times\mathbb P^1$ that is nef and not big, although its
fixed-density Gauduchon volume is exactly $n$.
A global projective coordinate function defines the background. An explicit
bounded radial potential compensates its determinant. Smooth radial
approximants give actual minimizing Gauduchon metrics for positive
perturbations, and the lower bound holds against every smooth positive
Gauduchon test. A distributional fiber average excludes every strict
current representative. The proof uses elementary calculations rather
than a Monge--Amp\`ere existence theorem or a nonclosed mass-invariance law.
\end{abstract}

We use $\ddc=\frac{i}{2\pi}\partial\bar\partial$ and ordinary wedge
powers. All quasi-psh potentials have their genuine canonical representatives.
Local integrability and positivity of current wedges with smooth forms are
used with their usual meanings; see \cite[Chapters I and III]{Dem}.
For a smooth strictly positive probability $\mu$ define
\[
 D_\mu(\gamma)=\int_X(\mu/\gamma^n)^{1/n}\gamma^n,
 \qquad\lambda_\mu(\theta)=\inf_{\ddc\gamma^{n-1}=0}
 \frac{\int_X\theta\wedge\gamma^{n-1}}{D_\mu(\gamma)},
 \qquad V_\mu(\theta)=\lambda_\mu(\theta)^n.
\]
The infimum includes every smooth positive Gauduchon metric $\gamma$.
Nefness means smooth strict potential representatives of $\theta+\epsilon h$
exist for every $\epsilon>0$, for a fixed Hermitian $h$.
Bigness means $\theta+\ddc v\geq\delta h$ for a global quasi-psh $v$
and $\delta>0$. These notions do not require closedness.

\begin{theorem}\label{th:main}
Let $m=n-1\geq2$, $t=(m-1)/m$, and let $\omega,w$ be normalized
Fubini--Study forms on $\mathbb P^m,\mathbb P^1$ respectively, so
$\int\omega^m=\int w=1$.
On $B=\mathbb P^m$ put
\[
 x([Z_0:\cdots:Z_m])=
 \frac{\sum_{j=1}^m|Z_j|^2}{\sum_{j=0}^m|Z_j|^2},
 \qquad a=x/t.
\]
On $X=B\times\mathbb P^1$, with pullbacks implicit, set
$h=\omega+w$, $\mu=\omega^m\wedge w$ and $\theta=\omega+a w$.
Then $\theta$ is smooth, semipositive, genuinely nonclosed, nef and not big,
while
\[
 \lambda_\mu(\theta)=n^{1/n},\qquad V_\mu(\theta)=n.
\]
For every $\epsilon>0$ the strictly positive background
$\theta_\epsilon=\omega+(a+\epsilon)w$ has
\[
 V_\mu(\theta_\epsilon)=\frac{n}{H_\epsilon},\qquad
 H_\epsilon=m\int_0^1\frac{y^{m-1}}{y/t+\epsilon}\,dy,
\]
and the proof specifies actual smooth positive Gauduchon minimizing metrics.
Hence $\theta$ disproves the unrestricted implication
$V_\mu(\theta)>0\Rightarrow\theta$ big for nef nonclosed backgrounds.
\end{theorem}

\section{Global smooth radial positive forms}
The function $x$ is smooth globally, ranges in $[0,1]$, vanishes only at
$p=[1:0:\cdots:0]$, and equals one on the hyperplane $Z_0=0$.
In the chart $Z_0=1$, write $s=|z|^2$; then $x=s/(1+s)$ and
$\omega=\ddc\log(1+s)$.
Polar integration of the Fubini--Study determinant $(1+s)^{-m-1}$,
with normalization $\int\omega^m=1$, gives
\begin{equation}\label{eq:radial}
 \int_B b(x)\omega^m=m\int_0^1 b(y)y^{m-1}\,dy.
\end{equation}
Indeed the radial density in $s$ is
$m s^{m-1}(1+s)^{-m-1}ds$, which becomes $m x^{m-1}dx$.
In particular $\int_B\omega^m/a=mt/(m-1)=1$.

For $\epsilon>0$ define $H_\epsilon$ as in the theorem and
\begin{equation}\label{eq:flux}
 q_\epsilon(x)^m=\frac{m}{H_\epsilon}
 \int_0^x\frac{y^{m-1}}{y/t+\epsilon}\,dy,
 \qquad
 \varphi_\epsilon(x)=\int_0^x
 \frac{q_\epsilon(y)-y}{y(1-y)}\,dy.
\end{equation}
The positive real root is used. At zero,
\[
 q_\epsilon(x)=x\left(\frac{m}{H_\epsilon}
 \int_0^1\frac{u^{m-1}}{xu/t+\epsilon}\,du\right)^{1/m},
\]
so it is $x$ times a smooth positive function. At one,
$q_\epsilon(1)=1$. Both apparent singularities in
$\varphi_\epsilon'$ are removable smoothly. Thus
$\varphi_\epsilon(x)$ is a global smooth base function.

For any radial local potential $F(s)$, put $q=sF'(s)$.
The tangential Hessian eigenvalue is $F'=q/s$ and the radial one is $dq/ds$.
For $F_\epsilon=\log(1+s)+\varphi_\epsilon(s/(1+s))$, direct
differentiation of \eqref{eq:flux} gives $sF_\epsilon'=q_\epsilon(x)$.
Both eigenvalues are positive, including at $s=0$ by the preceding smooth
factorization. If $A_\epsilon=\omega+\ddc\varphi_\epsilon$, the determinant
ratio in the affine chart is
\begin{equation}\label{eq:density}
 \frac{A_\epsilon^m}{\omega^m}
 =\frac{q_\epsilon^{m-1}\,dq_\epsilon/dx}{x^{m-1}}
 =\frac1{H_\epsilon(a+\epsilon)}.
\end{equation}
At the hyperplane $x=1$ the globally smooth form is the continuous limit
of positive forms, and its determinant ratio in \eqref{eq:density} tends
to a strictly positive number. Thus it is strictly positive there too.
Equivalently its tangential coefficient is that of $\omega$, and its normal
coefficient is $(dq_\epsilon/dx)(1)=1/[H_\epsilon(1/t+\epsilon)]>0$.
This proves global smooth strict positivity. It is closed by construction.

\section{A genuine bounded limiting potential}
By monotone convergence $H_\epsilon\to1$ as $\epsilon\downarrow0$.
For each $0<x\leq1$, \eqref{eq:flux} then gives
$q_\epsilon(x)\to x^t$.
We justify the limit of the potentials at the singular point as well.
Relative to the probability density $m y^{m-1}dy$, the weight
$(y/t+\epsilon)^{-1}$ is decreasing. Its normalized cumulative mass
is therefore at least the unweighted cumulative mass $x^m$.
Relative to the endpoint probability density $(m-1)y^{m-2}dy$, its weight
ratio is proportional to $y/(y+t\epsilon)$, which is increasing; its
normalized cumulative mass is at most $x^{m-1}$.
These two elementary comparisons follow by splitting either probability
at $x$: a decreasing weight has at least as large an average on $[0,x]$
as on $[x,1]$, and an increasing weight has the reverse inequality.
Taking positive $m$th roots yields
\begin{equation}\label{eq:bounds}
 x\leq q_\epsilon(x)\leq x^t.
\end{equation}

Define
\[
 \varphi(x)=\int_0^x\frac{y^t-y}{y(1-y)}\,dy.
\]
Its derivative is integrable at zero because $t>0$, and extends smoothly
at one with value $1-t$. Hence $\varphi$ is globally bounded continuous,
smooth off $p$, and $\varphi(p)=0$.
Equation~\eqref{eq:bounds} supplies the integrable derivative majorant
$(y^t-y)/[y(1-y)]$. Dominated convergence in $L^1([0,1])$ gives
\[
 \sup_{[0,1]}|\varphi_\epsilon-\varphi|\longrightarrow0.
\]
Consequently $A=\omega+\ddc\varphi$ is a positive closed current by
distributional convergence. Local smooth potentials for $\omega$ plus
$\varphi$ are continuous uniform limits of psh functions, so they are
genuinely psh, with the canonical value at $p$. This establishes admissibility
on the whole compact base, not only on its punctured affine chart.

Off $p$ the limiting radial flux is $q=x^t$, and the determinant formula
gives
\begin{equation}\label{eq:enddensity}
 \frac{A^m}{\omega^m}=
 \frac{(x^t)^{m-1}t x^{t-1}}{x^{m-1}}=t/x=1/a.
\end{equation}
The form is smooth strictly positive there, including at infinity by
the same determinant and continuity argument as before.
No power of this current at $p$ is needed in the following all-test bound.

\section{Exact volume over every smooth Gauduchon test}
On $X$, the probability $\mu$ is smooth strictly positive and
$\theta=\omega+a w$ is smooth semipositive. Thus $\theta+\eta h>0$
for every $\eta>0$, with zero perturbation potential, proving nefness.
Also $d\theta=da\wedge w\ne0$.
The bounded canonical potential pulled back from $B$ gives the positive
current $\Omega=\theta+\ddc\varphi=A+a w$. Off
$Z=\{p\}\times\mathbb P^1$ it is smooth strictly positive, and
\[
 \Omega^n=n a A^m\wedge w=n\mu.
\]
For every smooth positive Gauduchon $\gamma$, put $S=\gamma^{n-1}$.
Distributional Stokes gives $\int_X\theta\wedge S=\int_X\Omega\wedge S$,
because $\ddc S=0$ and $\varphi$ is bounded.
Matrix AM--GM off $Z$ and positivity of the global left measure imply
\begin{equation}\label{eq:lower}
 \int_X\theta\wedge\gamma^{n-1}\geq
 n^{1/n}\int_X(\mu/\gamma^n)^{1/n}\gamma^n.
\end{equation}
The set $Z$ has zero smooth volume; any singular left measure there only
strengthens this inequality. Thus $\lambda_\mu(\theta)\geq n^{1/n}$
over the entire original family. Using the closed metric $h$ shows finiteness.

We next give exact smooth minimizing tests for the positive perturbations.
Set $c_\epsilon=n/H_\epsilon$ and
\[
 \Omega_\epsilon=A_\epsilon+(a+\epsilon)w,
 \qquad\gamma_\epsilon=(a+\epsilon)^{-1/m}\Omega_\epsilon.
\]
These are smooth strictly positive. Equation~\eqref{eq:density} gives
$\Omega_\epsilon^n=c_\epsilon\mu$, and
\[
 T_\epsilon=\gamma_\epsilon^m
 =(a+\epsilon)^{-1}A_\epsilon^m
       +m A_\epsilon^{m-1}\wedge w
\]
is closed: its first term is a base top-degree form and its second term is
a wedge of closed forms. Therefore $\gamma_\epsilon$ is admissible.
Smooth Stokes for $\varphi_\epsilon$ gives exactly
\begin{align*}
 \int_X\theta_\epsilon\wedge T_\epsilon
 &=\int_X\Omega_\epsilon\wedge T_\epsilon
 =c_\epsilon\int_X(a+\epsilon)^{-1}\mu=n,\\
 D_\mu(\gamma_\epsilon)&=c_\epsilon^{m/n}H_\epsilon.
\end{align*}
Their quotient is $c_\epsilon^{1/n}$. For any smooth Gauduchon test,
the smooth positive representative $\Omega_\epsilon$ and matrix AM--GM
give the matching lower bound $c_\epsilon^{1/n}$. Hence
$V_\mu(\theta_\epsilon)=n/H_\epsilon$.
Since $\theta\leq\theta_\epsilon$, the same tests have unperturbed
quotient at most $c_\epsilon^{1/n}$. Letting $\epsilon\downarrow0$
and using \eqref{eq:lower} proves $V_\mu(\theta)=n$.

The decreasing nef definition with $\theta+\eta h$ agrees as well.
Its infima are at least $\lambda_\mu(\theta)$; for every fixed admissible
$\gamma$ its perturbed quotient converges to the unperturbed quotient.
Taking the infimum in this fixed-test upper bound proves the reverse limit
inequality. No smooth minimizer for the endpoint itself is asserted.

\section{Failure of bigness against all quasi-psh potentials}
Suppose $\theta+\ddc v\geq\delta h$ for a global genuine quasi-psh $v$
and $\delta>0$. Average over the holomorphic transitive fiber action
of $G=\mathrm{SU}(2)$ with Haar probability. The forms $\theta,h$ and
probability $\mu$ are invariant. Since $v\in L^1(\mu)$, Fubini defines
its average as an $L^1$ distribution; differentiation commutes with averaging
and the positive-current inequality is preserved.
The average is $q\circ\operatorname{pr}_1$ almost everywhere, where
$q\in L^1(B)$ is the fiber integral of $v$. Therefore
\[
 \theta+\ddc(q\circ\operatorname{pr}_1)-\delta h\geq0.
\]
Wedge this positive current with the smooth positive base top form
$\omega^m$. All horizontal terms vanish by base dimension, including the
distributional $\ddc(q\circ\operatorname{pr}_1)$. The resulting measure is
\[
 (a-\delta)\mu\geq0.
\]
It would imply $a\geq\delta$ almost everywhere. But $a$ is smooth,
$a(p)=0$, and $a<\delta$ on a nonempty open neighborhood of $p$.
This contradiction excludes every strict current representative.
No singular-current restriction to a fiber or alteration of a canonical
potential was used. It proves the nonbigness assertion.

\begin{remark}
The counterexample concerns the fixed smooth-density volume over all
Gauduchon tests. It does not replace that invariant by bounded-potential
Monge--Amp\`ere mass or an infimum over densities. The manifold is projective,
but its background is genuinely nonclosed; no closed nef K\"ahler theorem
is contradicted. The condition $m\geq2$ is essential, since it makes
$t>0$ and the limiting potential bounded. No priority or classification
claim is made.
\end{remark}
\begin{thebibliography}{9}
\bibitem{Dem} J.-P. Demailly, \emph{Complex Analytic and Differential Geometry},
authored manuscript, 21 June 2012; Chapter I, 4.17 and 5.1--5.3,
pp.~35 and 39; Chapter III, 1.13--1.16, pp.~132--134.
\url{https://people.math.harvard.edu/~siu/2dmml.pdf}.
\end{thebibliography}
\end{document}
